% Agent documentation:
% Purpose: describe the 24-month expansion path from Nesodden to the primary
% suburban segment.
% Downstream use: phase timing, markets, and thresholds feed the control
% framework.
% Revision guardrails: keep phase timing aligned with Sections 05, 07b, and
% 07d. Do not introduce Asker as a separate phase unless the goals change.
\subsection{Langsiktig plan fra Nesodden til primærsegmentet}

Den langsiktige planen starter ved første kommersielle lansering på
Nesodden. Måned 0 betyr derfor første måned med aktive bestillinger på
Nesodden. Denne tolkningen gjør at planen henger sammen med markedsmålet om
drift i minst to nye markeder innen 24 måneder etter første lansering (jf.\
\ref{sec:malsetninger}).

Planen bør ha tre faser. Fase 1 er Nesodden, som er beskrevet i forrige
del. Fase 2 starter i Bærum, der Zipline tester konkurranse mot Foodora og
Wolt i primærsegmentet. Asker står i segmenteringen som del av samme segment,
men bør bare legges til som replisering hvis Bærum-dataene er gode nok. Fase
3 er konsolidering og valg av neste geografiske område.

\subsubsection{Fase 2, Bærum først (måned 6--18)}

Bærum er første konkrete marked i primærsegmentet. Markedet har mange
eneboliger, høy kjøpekraft og etablerte leveringsvaner
\parencite{ssb_kommunefakta_baerum, wolt2026_asker_sandvika}. Her tester
Zipline om tjenesten kan vinne neste ordre fra Foodora og Wolt.

Tiltakene i Bærum bør skille seg tydelig fra Nesodden-tiltakene. På Nesodden
skal Zipline bygge tillit til en ny leveringsform. I Bærum skal Zipline vinne
en ordre kunden ellers ville lagt hos Foodora eller Wolt. Derfor bør Zipline
bruke korte appkampanjer, partnerkampanjer, oppfølgingsannonsering i
leveringssonen og tydelig sanntidssporing etter kjøp.

CAC per aktiv kunde påvirkes først og fremst gjennom presis målretting.
Annonser bør avgrenses til husholdninger i dekningssonen og til digitale
flater der matlevering allerede er relevant. Da brukes ikke budsjettet på
kunder som ikke kan bestille. Konvertering til første kjøp bør måles fra
appbesøk i dekningssonen til første ordre. Appen bør ha en enkel flyt med
partner, varevalg, handlekurv, betaling og sporing, slik at dronelevering
ikke oppleves som et ekstra valg kunden må lære.

Pris- og tidsfordelen bør vises i selve kampanjen og i appen før kunden
bekrefter kjøpet. Budskapet bør vise lavere leveringsgebyr og kortere
leveringstid enn alternativene kunden allerede bruker. En byttekampanje med
gratis eller sterkt rabattert første levering kan løfte konverteringen, men
tilbudet bør tidsavgrenses slik at CAC ikke blir varig høy. Etter kjøp bør
Zipline spørre hva kunden ellers ville brukt. Slik måles andelen ordre som
erstatter Foodora, Wolt, Oda, egen henting eller ingen bestilling.

Snittkurv bør måles, men komme etter kundevekst og gjenkjøp. Platform~2 kan
frakte 3,6~kg. Derfor bør ordreverdien økes gjennom små tillegg og ferdige
pakker, ikke gjennom store handlekurver. Restauranter og butikker i Sandvika
og Lysaker kan bruke middagspakker, drikke, dessert og terskelrabatter når
kunden allerede har lagt noe i kurven. Målet er å øke ordreverdien uten å
gjøre kjøpet mer komplisert.

Ved måned 12 bør Zipline vurdere Bærum før Asker eventuelt legges til som
replisering innen samme segment. Ved måned 18 bør Bærum og eventuell
drift i Asker vurderes opp mot markedsmålet om minst 3 prosent adopsjonsrate
i nye markeder. Aktiv kunde bør fortsatt defineres som minst én bestilling
siste 30 dager. Gjenkjøp og CAC per aktiv kunde bør inngå i samme vurdering.

\subsubsection{Fase 3, konsolidering og neste marked (måned 18--24)}

I siste fase bør Zipline stabilisere Nesodden, Bærum og eventuell
drift i Asker før nye områder legges til. Hovedoppgaven er å sammenligne
aktiv kundebase, gjenkjøp, CAC per aktiv kunde, andel ordre som erstatter
substitutter, punktlighet, kundetilfredshet og partnerresultater på tvers av
markedene. Snittkurv og bidragsmargin bør måles, men ikke styre
markedsføringen alene.

Sammenligningen bør styre neste geografiske valg. Hvis Nesodden gir lav CAC
per aktiv kunde og høyt gjenkjøp, bør Zipline vurdere flere
logistikkisolerte områder rundt Oslofjorden. Hvis Bærum gir bedre
enhetsøkonomi, bør selskapet prioritere flere suburbane eneboligområder.
Beslutningen bør bygge på data, ikke bare på markedsstørrelse.

\subsubsection{Tverrgående arbeid}

Regulatorisk arbeid og partnerarbeid må gå gjennom hele planen.
Driftsgodkjenninger fra Luftfartstilsynet og EASA avgjør hvor raskt Zipline
kan skalere \parencite{luftfartstilsynet2023, euuas2021}. Partnerarbeidet
avgjør om tjenesten har varer kundene faktisk vil bestille.

Planen bygger på tre antagelser som må følges løpende. Først antas det at
Wing \parencite{wing2024} og andre droneleveringsselskaper ikke etablerer seg
tungt i Norge i perioden. Deretter antas det at Foodora og Wolt ikke raskt kan
kopiere dronelevering gjennom eget regelverksgjennombrudd. Til slutt antas det
at norske vinterforhold ikke gir så mange avbrudd at kundeopplevelsen svekkes.
Hvis en av antagelsene faller, må planen revideres.
